% ======================================================================
% fire_section.tex -- fire model, calibration, population, learning k_F
% % ======================================================================
% fire_section.tex -- fire model, calibration, population, learning k_F
% % ======================================================================
% fire_section.tex -- fire model, calibration, population, learning k_F
% % ======================================================================
% fire_section.tex -- fire model, calibration, population, learning k_F
% \input{fire_section} from main.tex. Needs: amsmath, siunitx, booktabs,
% cleveref, and the assumption environment defined in main.tex.
% Lines marked % UPDATE depend on the geometry and must be refreshed after
% rerunning learn.py, refine.py, calibrate_spread.py and learn_kf.py.
% ======================================================================

\subsection{Fire Scenario}\label{sec:fire}

We extend the evacuation model with a fire that starts at a random location in the
hall and spreads while patrons evacuate. The fire enters the model in three ways: it
makes cells impassable, it repels patrons through a fire proximity field, and it causes
casualties among patrons who remain in burning cells too long. Throughout, one time step
is $\Delta t = \SI{0.36}{s}$ and one cell is \SI{0.432}{m} square.

\subsubsection{Fire Spread}\label{sec:fire_spread}

We model fire with a stochastic cellular automaton adapted from the fire mechanics of
Minecraft (Java Edition) \cite{minecraftfire}, restricted to two dimensions. Each cell
$c$ carries an age $A_c \in \{-1, 0, 1, \dots, 15\}$, where $A_c = -1$ means $c$ is not
burning, and a flag $B_c \in \{0, 1\}$ recording whether its fuel has been used up. A cell
is \emph{flammable} if it is not a wall, not burning, and $B_c = 0$. Let $M(c)$ denote the
eight cells surrounding $c$.

At each time step the fire updates with probability $q$ (analogous to Minecraft's random
block ticks); when it updates, every cell changes simultaneously:
\begin{enumerate}
  \item \emph{Spread.} A burning cell $c'$ ignites each flammable neighbor independently
  with probability
  \begin{equation}\label{eq:fire_spread}
    p_{\text{ign}}(A_{c'}) = s \cdot \frac{\iota + 40 + 7\kappa}{100\,(A_{c'} + 30)},
  \end{equation}
  so a flammable cell $c$ catches fire with probability
  \begin{equation}\label{eq:fire_ignite}
    P_{\text{ign}}(c) = 1 - \prod_{\substack{c' \in M(c) \\ A_{c'} \ge 0}}
    \bigl(1 - p_{\text{ign}}(A_{c'})\bigr).
  \end{equation}
  \item \emph{Consumption.} A burning cell uses up its fuel with probability
  \begin{equation}\label{eq:fire_consume}
    p_{\text{cons}} = \frac{\beta}{300}.
  \end{equation}
  \item \emph{Aging.} The age of a burning cell increases by one with probability $1/3$,
  up to $15$.
  \item \emph{Burnout.} A burning cell goes out if $A_c > 3$ and no neighbor is flammable,
  or with probability $1/4$ once $A_c = 15$.
\end{enumerate}
A cell whose fire is consumed or burns out becomes ash: it is passable but no longer
flammable. Newly ignited cells start at age $0$. Each simulation begins with a single
ignition cell drawn uniformly at random from the non-wall, non-exit, unoccupied cells.

The functional forms of \cref{eq:fire_spread,eq:fire_consume} and the aging and burnout
rules follow Minecraft. Minecraft's fire is a block that sits next to separate fuel blocks
and can spread several blocks upward; in two dimensions we let the fire occupy the fuel
cell itself and spread within its eight-cell neighborhood. The constants are those of
wool in Minecraft, ignite odds $\iota = 30$ and burn odds $\beta = 60$, at the normal
difficulty level $\kappa = 2$ \cite{minecraftfire}, so a burning cell uses up its fuel
with probability $0.2$ per update. The spread scale $s$ and update probability $q$ are
calibrated in \cref{sec:fire_calibration}; Minecraft itself corresponds to $s = 1$.

\subsubsection{Calibrating the Fire}\label{sec:fire_calibration}

The spread scale $s$ and the update probability $q$ play different roles. The spread
scale controls the \emph{shape} of the fire: whether an ignition becomes an established
fire and how completely the fire front burns the hall. The update probability controls
only its \emph{speed}: lowering $q$ stretches the fire's timeline by a factor $1/q$
without changing how it spreads. We calibrate them separately.

\paragraph{Shape.} For each candidate $s$ we ran the fire alone from $40$ random ignition
sites, updating every step ($q = 1$). For each run we recorded the coverage $C(t)$, the
fraction of non-wall cells that have caught fire by update $t$, and its final value
$C_{\text{final}}$. We call a fire \emph{established} if $C_{\text{final}} \ge 0.5$, and
define its \emph{engulfment time} as the first $t$ with
$C(t) \ge 0.95\, C_{\text{final}}$. Because the front burns out behind itself and leaves
unburnt pockets, $C_{\text{final}} < 1$ even for fires that sweep the hall.
\Cref{tab:fire_calibration} shows the results. Below $s \approx 3$ almost no ignition
becomes established; the fire behaves like a branching process near its critical point,
so lowering $s$ makes fires fizzle rather than burn more slowly. We chose $s = 8$, for
which about $95\%$ of ignitions become established fires that burn essentially the whole
hall. This deliberately keeps a small share of fires that never take hold, as happens in
reality. Larger $s$ pushes $p_{\text{ign}}$ toward $1$, where the front becomes a
deterministic square determined by the grid rather than by the fire.

\begin{table}[H]
  \centering
  \caption{Fire calibration with the fire updating every step ($q = 1$), $40$ random
  ignitions per value of $s$. Coverage and engulfment updates are medians over
  established fires; $q$ is the update probability giving a median engulfment time of
  \SI{180}{s}.}
  \label{tab:fire_calibration}
  \begin{tabular}{rcccc}
    \toprule
    $s$ & established & coverage & engulfment (updates) & $q$ for \SI{180}{s} \\
    \midrule
    4  & $80\%$  & $0.90$ & $197$ & $0.394$ \\
    6  & $88\%$  & $0.99$ & $107$ & $0.214$ \\
    8  & $95\%$  & $1.00$ & $83$  & $0.166$ \\
    12 & $100\%$ & $1.00$ & $64$  & $0.128$ \\
    16 & $100\%$ & $1.00$ & $56$  & $0.113$ \\
    \bottomrule
  \end{tabular}
  % UPDATE: rerun calibrate_spread.py with the measured access points
\end{table}

\paragraph{Speed.} With $s = 8$ fixed, an established fire takes a median of $83$ updates
to engulf the hall, so an engulfment time $T_{\text{eng}}$ requires
\begin{equation}\label{eq:fire_q}
  q = \frac{83\,\Delta t}{T_{\text{eng}}} = 0.166 \cdot \frac{\SI{180}{s}}{T_{\text{eng}}}.
\end{equation}
We use $T_{\text{eng}} = \SI{6}{min}$, giving $q = 0.083$.

The engulfment time is the least certain input of the fire model, and we treat it as a
scenario parameter rather than a physical constant. For comparison, standard
fire-protection practice models fire growth as a $t^2$ design fire, with heat release
rate $Q(t) = \alpha t^2$; upholstered furniture is usually classed as a ``fast'' fire,
$\alpha = \SI{0.0469}{kW/s^2}$ \cite{sfpe}. If the burning area is $Q/q''$ for a heat
release rate per unit area $q''$, the flame front advances at
$v_f = \sqrt{\alpha/(\pi q'')}$, only a few millimeters per second for upholstered
seating. Flames that slow would not reach most patrons during an evacuation lasting about
a minute. In real assembly fires, smoke rather than flame spread is the main hazard, and
it fills a hall far faster than flames spread \cite{sfpe}. Our engulfment time should
therefore be read as the time for conditions throughout the hall to become untenable,
with the fire field standing in for flames and smoke together. We first used
$T_{\text{eng}} = \SI{3}{min}$, which produced $13$--$16\%$ casualties; we regard that as
implausibly severe and use \SI{6}{min}, reporting the \SI{3}{min} case as a sensitivity
check.

\subsubsection{Fire Proximity Field}\label{sec:fire_field}

Patrons should react to the fire before they reach it. Let $r_{ij}$ be the number of steps
from $(i, j)$ to the nearest burning cell through non-wall cells, using the eight-cell
neighborhood. The fire proximity field is
\begin{equation}\label{eq:fire_field}
  F_{ij} = \max\bigl(0,\ \ell - r_{ij}\bigr), \qquad \ell = 8 \text{ cells}
  \ (\approx \SI{3.5}{m}).
\end{equation}
Within $\ell$ cells of the fire, each step away from it lowers $F$ by exactly one, which is
the same unit in which each step toward an exit lowers $SD$. A fire weight $k_F$ is then
directly comparable to the exit weight $k_S$: $k_F / k_S$ is how much a patron values one
step away from the fire relative to one step toward an exit. An exponentially decaying
field would make $k_F$ depend on the distance at which it is evaluated and lose this
interpretation. Beyond $\ell$ the fire has no direct influence on movement.

\subsubsection{Movement and Casualties}\label{sec:fire_movement}

With fire present, the transition probability of patron $a$ at $(i, j)$ is
\begin{equation}\label{eq:trans_prob_fire}
  p^{(a)}_{i',j'} = N^{(a)}_{i,j}
  \exp\!\Bigl(k^{(a)}_D D_{i',j'} - k^{(a)}_S\, SD_{i',j'} - k^{(a)}_F F_{i',j'}\Bigr)
  \,\mathbf{1}\bigl[(i', j') \text{ admissible}\bigr],
\end{equation}
over the eight neighboring cells and the patron's own cell. A target cell is admissible if
it is walkable, unoccupied (by a patron or a casualty), and not burning. Because burning
cells and casualties block paths, $SD$ is recomputed whenever either changes. Cells from
which no exit can be reached are assigned one more than the largest finite distance, so
patrons cut off from every exit still move away from danger.

A patron standing in a burning cell accumulates one step of exposure per time step. When a
patron's cumulative exposure reaches \SI{5}{s} ($14$ steps), the patron becomes a
casualty and remains in the cell as an obstacle for the rest of the simulation; until
then, the patron can still walk out of the fire. Patrons still inside at the time horizon
$T = 2000$ steps (\SI{12}{min}) are counted as trapped. Both burned and trapped patrons
count as casualties.

\subsubsection{Heterogeneous Population}\label{sec:population}

Each patron $a$ receives individual weights $(k^{(a)}_S, k^{(a)}_D, k^{(a)}_F)$. By
\cref{as:optimal}, patrons behave approximately optimally under normal conditions, so we
draw $(k_S, k_D)$ from a neighborhood of the optimum $x^* = (k_S^*, k_D^*)$ for a
homogeneous crowd without fire. Near that optimum the mean exit time is well approximated
by a quadratic,
\begin{equation}\label{eq:quadratic}
  J(x) \approx J^* + \tfrac{1}{2}(x - x^*)^\top H\, (x - x^*),
\end{equation}
which we fit by least squares to the refined grid of simulations, obtaining $x^*$ and the
Hessian $H$. We draw
\begin{equation}\label{eq:population}
  \bigl(k^{(a)}_S, k^{(a)}_D\bigr) \sim \mathcal{N}\bigl(x^*,\ c\,H^{-1}\bigr),
\end{equation}
redrawing any sample with $k_S \le 0$. Two properties motivate this choice. First, patrons
vary widely along directions in which $J$ is flat and narrowly along directions in which
it is steep. Second, the spread has a direct meaning: if each patron's parameters were
adopted by a whole crowd,
\begin{equation}\label{eq:expected_loss}
  \mathbb{E}\bigl[J(x^{(a)}) - J^*\bigr]
  \approx \tfrac{1}{2}\operatorname{tr}\bigl(H \cdot c\,H^{-1}\bigr) = c,
\end{equation}
so $c$ is the average number of time steps lost per patron. Equivalently,
\cref{eq:population} is the quadratic approximation of a Boltzmann distribution
$p(x) \propto \exp\bigl(-(J(x) - J^*)/c\bigr)$, in which cheaper behaviors are more common
and $c$ acts as a temperature. We use $c = 1$. With the fitted
$x^* = (1.89, -3.02)$ and $H$, this gives % UPDATE
$\operatorname{sd}(k_S) \approx 0.35$, $\operatorname{sd}(k_D) \approx 2.45$ and
$\operatorname{corr}(k_S, k_D) \approx -0.76$: patrons agree closely on how strongly they
head for an exit but differ widely in how they respond to the crowd. % UPDATE

The fire weight must be positive, so we draw it from a lognormal distribution,
independently of $(k_S, k_D)$:
\begin{equation}\label{eq:kF}
  k^{(a)}_F = m\, e^{\sigma_F z_a}, \qquad z_a \sim \mathcal{N}(0, 1),
\end{equation}
with median $m$ and $\sigma_F = 0.5$, so about two thirds of patrons have $k_F$ within a
factor $e^{0.5} \approx 1.65$ of $m$. The median $m$ is learned in \cref{sec:learn_kF}.

\subsubsection{Learning the Fire Weight}\label{sec:learn_kF}

Consistent with \cref{as:consistent}, we hold the distribution of $(k_S, k_D)$ fixed and
learn only the median fire weight $m$. For one simulation, let $C$ be the casualty
fraction and $\bar t$ the mean exit time of survivors. We minimize
\begin{equation}\label{eq:fire_objective}
  \mathcal{J}(m) = \mathbb{E}\left[\, C + \lambda\, \frac{\bar t}{t_{\text{ref}}} \right],
  \qquad \lambda = 0.1,\ t_{\text{ref}} = 150 \text{ steps},
\end{equation}
where the expectation is over the ignition site, the population draw, the fire, and the
patrons' random choices; $t_{\text{ref}}$ is the optimal mean exit time without fire. With
these values one percent of casualties counts as much as $15$ steps of survivor exit time,
so casualties are the primary objective and exit time breaks near-ties.

We estimated $\mathcal{J}(m)$ for $m \in \{0, 0.5, 1, 2, 3, 5, 8, 12, 20\}$ with $R = 50$
replicates per value. Each replicate fixes an ignition site, a population draw (including
the standard normals $z_a$ in \cref{eq:kF}), a random stream for the patrons and a separate
random stream for the fire. These are shared across all values of $m$ (common random
numbers), and since the fire does not depend on the patrons, each replicate's fire evolves
identically for every $m$. Differences between values of $m$ are therefore paired
comparisons on identical scenarios.

\paragraph{Results.} \Cref{tab:kF} shows the results. % UPDATE all numbers below
Moderate fire avoidance roughly halves casualties: raising the median fire weight from
$m = 0$ to $m = 1$ lowers the casualty fraction from $2.3\%$ to $1.3\%$ (about $13$ of
$1010$ patrons). Strong avoidance is counterproductive: casualties rise to $3.1\%$ at
$m = 5$ and $3.8\%$ at $m = 20$, about three times the optimum, and the share of patrons
trapped with no route to an exit grows with $m$. Patrons who put distance from the fire
ahead of progress toward an exit are cut off or caught in worse positions. Fire avoidance
also delays survivors, whose mean exit time rises from $151$ steps at $m = 0$ to $180$ at
$m = 20$. $\mathcal{J}$ is minimized at $m = 1$, and no other value is within two paired
standard errors of it. Relative to $k_S^* \approx 1.9$, the best response weights a step
away from the fire about half as much as a step toward an exit.

\begin{table}[H]
  \centering
  \caption{Learning the median fire weight $m$ ($T_{\text{eng}} = \SI{6}{min}$, $50$
  replicates). Means over replicates; casualties are burned plus trapped, exit time is the
  survivors' mean in steps.}
  \label{tab:kF}
  % UPDATE: rerun learn_kf.py with the measured access points
  \begin{tabular}{rccccc}
    \toprule
    $m$ & $\mathcal{J}$ & casualties & burned & trapped & exit time \\
    \midrule
    0    & $0.124$ & $2.29\%$ & $1.93\%$ & $0.36\%$ & $151.1$ \\
    0.5  & $0.119$ & $1.66\%$ & $1.48\%$ & $0.18\%$ & $153.3$ \\
    1    & $0.116$ & $1.29\%$ & $1.08\%$ & $0.21\%$ & $155.0$ \\
    2    & $0.127$ & $1.87\%$ & $1.52\%$ & $0.35\%$ & $162.0$ \\
    3    & $0.135$ & $2.30\%$ & $1.88\%$ & $0.43\%$ & $168.0$ \\
    5    & $0.148$ & $3.09\%$ & $2.55\%$ & $0.54\%$ & $175.2$ \\
    8    & $0.153$ & $3.43\%$ & $2.84\%$ & $0.59\%$ & $177.7$ \\
    12   & $0.157$ & $3.79\%$ & $3.14\%$ & $0.65\%$ & $179.0$ \\
    20   & $0.159$ & $3.84\%$ & $3.16\%$ & $0.68\%$ & $180.1$ \\
    \bottomrule
  \end{tabular}
\end{table}

\paragraph{Sensitivity.} The optimum is robust to the fire model. With a faster fire
($T_{\text{eng}} = \SI{3}{min}$) and casualties occurring on first contact with the fire,
casualties are far higher ($13$--$16\%$), but $\mathcal{J}$ is again minimized at $m = 1$.
% UPDATE if rerun

\subsubsection{Constants}\label{sec:fire_constants}

\Cref{tab:fire_constants} lists every constant in the fire model and how it was determined.

\begin{table}[H]
  \centering
  \caption{Constants of the fire model.}
  \label{tab:fire_constants}
  \small
  \begin{tabular}{llll}
    \toprule
    Symbol & Meaning & Value & Source \\
    \midrule
    $\Delta t$ & time step & \SI{0.36}{s} & walking speed, cell size \\
    $\iota$ & ignite odds of fuel & $30$ & Minecraft (wool) \\
    $\beta$ & burn odds of fuel & $60$ & Minecraft (wool) \\
    $\kappa$ & difficulty level & $2$ & Minecraft (normal) \\
    $s$ & spread scale & $8$ & calibrated: $\approx 95\%$ of ignitions establish \\
    $T_{\text{eng}}$ & median engulfment time & \SI{6}{min} & assumed; \SI{3}{min} as sensitivity \\
    $q$ & fire update probability & $0.083$ & calibrated from $s$ and $T_{\text{eng}}$, \cref{eq:fire_q} \\
    $\ell$ & fire sensing range & $8$ cells (\SI{3.5}{m}) & assumed \\
    -- & exposure before casualty & \SI{5}{s} ($14$ steps) & assumed \\
    $T$ & time horizon & $2000$ steps & longer than any evacuation \\
    $x^*$ & center of $(k_S, k_D)$ & $(1.89, -3.02)$ & learned without fire % UPDATE
    \\
    $H$ & curvature of $J$ at $x^*$ & -- & fitted, \cref{eq:quadratic} \\
    $c$ & population spread & $1$ step & assumed \\
    $m$ & median fire weight & $1$ & learned, \cref{sec:learn_kF} % UPDATE
    \\
    $\sigma_F$ & spread of $k_F$ & $0.5$ & assumed \\
    $\lambda,\ t_{\text{ref}}$ & objective weights & $0.1,\ 150$ steps & chosen, \cref{eq:fire_objective} \\
    \bottomrule
  \end{tabular}
\end{table}

\subsubsection{Fire Assumptions}\label{sec:fire_assumptions}

\begin{assumption}
Every non-wall cell (seats, legroom, aisles, access points and exits) is equally
flammable, with the properties of wool in Minecraft; walls are stone and never burn.
\end{assumption}
\begin{assumption}
Fire spreads only between adjacent cells, and its spread does not depend on the patrons.
\end{assumption}
\begin{assumption}
The fire field represents untenable conditions (flames and smoke together), which spread
through the hall in a median of \SI{6}{min}.
\end{assumption}
\begin{assumption}
Patrons sense the fire within \SI{3.5}{m} and react to it in proportion to how many steps
away from it a move takes them.
\end{assumption}
\begin{assumption}
A patron who spends \SI{5}{s} in total in burning cells is incapacitated and becomes an
obstacle to others.
\end{assumption}
\begin{assumption}
Patrons' exit-seeking and crowd behavior, $(k_S, k_D)$, are the same as without fire
(\cref{as:consistent}); only their response to the fire, $k_F$, is new.
\end{assumption}

\subsubsection{Limitations}\label{sec:fire_limitations}

The fire model is two-dimensional and does not represent smoke, heat or the hot gas layer
under the ceiling, which cause most deaths in real assembly fires; the engulfment time is
a single stylized timescale for all of these. The Minecraft spread rules are not derived
from combustion physics, and their constants follow the game rather than measurements of
theater seating. Patrons respond only to fire within \SI{3.5}{m}; an alarm, the sight of
smoke across the hall or staff instructions would alert them sooner. Casualties are
binary, with no injuries or reduced walking speed. Finally, the optimal $(k_S, k_D)$ were
learned for a homogeneous crowd without fire, so the population in
\cref{eq:population} is near-optimal in that sense, not necessarily under fire.
 from main.tex. Needs: amsmath, siunitx, booktabs,
% cleveref, and the assumption environment defined in main.tex.
% Lines marked % UPDATE depend on the geometry and must be refreshed after
% rerunning learn.py, refine.py, calibrate_spread.py and learn_kf.py.
% ======================================================================

\subsection{Fire Scenario}\label{sec:fire}

We extend the evacuation model with a fire that starts at a random location in the
hall and spreads while patrons evacuate. The fire enters the model in three ways: it
makes cells impassable, it repels patrons through a fire proximity field, and it causes
casualties among patrons who remain in burning cells too long. Throughout, one time step
is $\Delta t = \SI{0.36}{s}$ and one cell is \SI{0.432}{m} square.

\subsubsection{Fire Spread}\label{sec:fire_spread}

We model fire with a stochastic cellular automaton adapted from the fire mechanics of
Minecraft (Java Edition) \cite{minecraftfire}, restricted to two dimensions. Each cell
$c$ carries an age $A_c \in \{-1, 0, 1, \dots, 15\}$, where $A_c = -1$ means $c$ is not
burning, and a flag $B_c \in \{0, 1\}$ recording whether its fuel has been used up. A cell
is \emph{flammable} if it is not a wall, not burning, and $B_c = 0$. Let $M(c)$ denote the
eight cells surrounding $c$.

At each time step the fire updates with probability $q$ (analogous to Minecraft's random
block ticks); when it updates, every cell changes simultaneously:
\begin{enumerate}
  \item \emph{Spread.} A burning cell $c'$ ignites each flammable neighbor independently
  with probability
  \begin{equation}\label{eq:fire_spread}
    p_{\text{ign}}(A_{c'}) = s \cdot \frac{\iota + 40 + 7\kappa}{100\,(A_{c'} + 30)},
  \end{equation}
  so a flammable cell $c$ catches fire with probability
  \begin{equation}\label{eq:fire_ignite}
    P_{\text{ign}}(c) = 1 - \prod_{\substack{c' \in M(c) \\ A_{c'} \ge 0}}
    \bigl(1 - p_{\text{ign}}(A_{c'})\bigr).
  \end{equation}
  \item \emph{Consumption.} A burning cell uses up its fuel with probability
  \begin{equation}\label{eq:fire_consume}
    p_{\text{cons}} = \frac{\beta}{300}.
  \end{equation}
  \item \emph{Aging.} The age of a burning cell increases by one with probability $1/3$,
  up to $15$.
  \item \emph{Burnout.} A burning cell goes out if $A_c > 3$ and no neighbor is flammable,
  or with probability $1/4$ once $A_c = 15$.
\end{enumerate}
A cell whose fire is consumed or burns out becomes ash: it is passable but no longer
flammable. Newly ignited cells start at age $0$. Each simulation begins with a single
ignition cell drawn uniformly at random from the non-wall, non-exit, unoccupied cells.

The functional forms of \cref{eq:fire_spread,eq:fire_consume} and the aging and burnout
rules follow Minecraft. Minecraft's fire is a block that sits next to separate fuel blocks
and can spread several blocks upward; in two dimensions we let the fire occupy the fuel
cell itself and spread within its eight-cell neighborhood. The constants are those of
wool in Minecraft, ignite odds $\iota = 30$ and burn odds $\beta = 60$, at the normal
difficulty level $\kappa = 2$ \cite{minecraftfire}, so a burning cell uses up its fuel
with probability $0.2$ per update. The spread scale $s$ and update probability $q$ are
calibrated in \cref{sec:fire_calibration}; Minecraft itself corresponds to $s = 1$.

\subsubsection{Calibrating the Fire}\label{sec:fire_calibration}

The spread scale $s$ and the update probability $q$ play different roles. The spread
scale controls the \emph{shape} of the fire: whether an ignition becomes an established
fire and how completely the fire front burns the hall. The update probability controls
only its \emph{speed}: lowering $q$ stretches the fire's timeline by a factor $1/q$
without changing how it spreads. We calibrate them separately.

\paragraph{Shape.} For each candidate $s$ we ran the fire alone from $40$ random ignition
sites, updating every step ($q = 1$). For each run we recorded the coverage $C(t)$, the
fraction of non-wall cells that have caught fire by update $t$, and its final value
$C_{\text{final}}$. We call a fire \emph{established} if $C_{\text{final}} \ge 0.5$, and
define its \emph{engulfment time} as the first $t$ with
$C(t) \ge 0.95\, C_{\text{final}}$. Because the front burns out behind itself and leaves
unburnt pockets, $C_{\text{final}} < 1$ even for fires that sweep the hall.
\Cref{tab:fire_calibration} shows the results. Below $s \approx 3$ almost no ignition
becomes established; the fire behaves like a branching process near its critical point,
so lowering $s$ makes fires fizzle rather than burn more slowly. We chose $s = 8$, for
which about $95\%$ of ignitions become established fires that burn essentially the whole
hall. This deliberately keeps a small share of fires that never take hold, as happens in
reality. Larger $s$ pushes $p_{\text{ign}}$ toward $1$, where the front becomes a
deterministic square determined by the grid rather than by the fire.

\begin{table}[H]
  \centering
  \caption{Fire calibration with the fire updating every step ($q = 1$), $40$ random
  ignitions per value of $s$. Coverage and engulfment updates are medians over
  established fires; $q$ is the update probability giving a median engulfment time of
  \SI{180}{s}.}
  \label{tab:fire_calibration}
  \begin{tabular}{rcccc}
    \toprule
    $s$ & established & coverage & engulfment (updates) & $q$ for \SI{180}{s} \\
    \midrule
    4  & $80\%$  & $0.90$ & $197$ & $0.394$ \\
    6  & $88\%$  & $0.99$ & $107$ & $0.214$ \\
    8  & $95\%$  & $1.00$ & $83$  & $0.166$ \\
    12 & $100\%$ & $1.00$ & $64$  & $0.128$ \\
    16 & $100\%$ & $1.00$ & $56$  & $0.113$ \\
    \bottomrule
  \end{tabular}
  % UPDATE: rerun calibrate_spread.py with the measured access points
\end{table}

\paragraph{Speed.} With $s = 8$ fixed, an established fire takes a median of $83$ updates
to engulf the hall, so an engulfment time $T_{\text{eng}}$ requires
\begin{equation}\label{eq:fire_q}
  q = \frac{83\,\Delta t}{T_{\text{eng}}} = 0.166 \cdot \frac{\SI{180}{s}}{T_{\text{eng}}}.
\end{equation}
We use $T_{\text{eng}} = \SI{6}{min}$, giving $q = 0.083$.

The engulfment time is the least certain input of the fire model, and we treat it as a
scenario parameter rather than a physical constant. For comparison, standard
fire-protection practice models fire growth as a $t^2$ design fire, with heat release
rate $Q(t) = \alpha t^2$; upholstered furniture is usually classed as a ``fast'' fire,
$\alpha = \SI{0.0469}{kW/s^2}$ \cite{sfpe}. If the burning area is $Q/q''$ for a heat
release rate per unit area $q''$, the flame front advances at
$v_f = \sqrt{\alpha/(\pi q'')}$, only a few millimeters per second for upholstered
seating. Flames that slow would not reach most patrons during an evacuation lasting about
a minute. In real assembly fires, smoke rather than flame spread is the main hazard, and
it fills a hall far faster than flames spread \cite{sfpe}. Our engulfment time should
therefore be read as the time for conditions throughout the hall to become untenable,
with the fire field standing in for flames and smoke together. We first used
$T_{\text{eng}} = \SI{3}{min}$, which produced $13$--$16\%$ casualties; we regard that as
implausibly severe and use \SI{6}{min}, reporting the \SI{3}{min} case as a sensitivity
check.

\subsubsection{Fire Proximity Field}\label{sec:fire_field}

Patrons should react to the fire before they reach it. Let $r_{ij}$ be the number of steps
from $(i, j)$ to the nearest burning cell through non-wall cells, using the eight-cell
neighborhood. The fire proximity field is
\begin{equation}\label{eq:fire_field}
  F_{ij} = \max\bigl(0,\ \ell - r_{ij}\bigr), \qquad \ell = 8 \text{ cells}
  \ (\approx \SI{3.5}{m}).
\end{equation}
Within $\ell$ cells of the fire, each step away from it lowers $F$ by exactly one, which is
the same unit in which each step toward an exit lowers $SD$. A fire weight $k_F$ is then
directly comparable to the exit weight $k_S$: $k_F / k_S$ is how much a patron values one
step away from the fire relative to one step toward an exit. An exponentially decaying
field would make $k_F$ depend on the distance at which it is evaluated and lose this
interpretation. Beyond $\ell$ the fire has no direct influence on movement.

\subsubsection{Movement and Casualties}\label{sec:fire_movement}

With fire present, the transition probability of patron $a$ at $(i, j)$ is
\begin{equation}\label{eq:trans_prob_fire}
  p^{(a)}_{i',j'} = N^{(a)}_{i,j}
  \exp\!\Bigl(k^{(a)}_D D_{i',j'} - k^{(a)}_S\, SD_{i',j'} - k^{(a)}_F F_{i',j'}\Bigr)
  \,\mathbf{1}\bigl[(i', j') \text{ admissible}\bigr],
\end{equation}
over the eight neighboring cells and the patron's own cell. A target cell is admissible if
it is walkable, unoccupied (by a patron or a casualty), and not burning. Because burning
cells and casualties block paths, $SD$ is recomputed whenever either changes. Cells from
which no exit can be reached are assigned one more than the largest finite distance, so
patrons cut off from every exit still move away from danger.

A patron standing in a burning cell accumulates one step of exposure per time step. When a
patron's cumulative exposure reaches \SI{5}{s} ($14$ steps), the patron becomes a
casualty and remains in the cell as an obstacle for the rest of the simulation; until
then, the patron can still walk out of the fire. Patrons still inside at the time horizon
$T = 2000$ steps (\SI{12}{min}) are counted as trapped. Both burned and trapped patrons
count as casualties.

\subsubsection{Heterogeneous Population}\label{sec:population}

Each patron $a$ receives individual weights $(k^{(a)}_S, k^{(a)}_D, k^{(a)}_F)$. By
\cref{as:optimal}, patrons behave approximately optimally under normal conditions, so we
draw $(k_S, k_D)$ from a neighborhood of the optimum $x^* = (k_S^*, k_D^*)$ for a
homogeneous crowd without fire. Near that optimum the mean exit time is well approximated
by a quadratic,
\begin{equation}\label{eq:quadratic}
  J(x) \approx J^* + \tfrac{1}{2}(x - x^*)^\top H\, (x - x^*),
\end{equation}
which we fit by least squares to the refined grid of simulations, obtaining $x^*$ and the
Hessian $H$. We draw
\begin{equation}\label{eq:population}
  \bigl(k^{(a)}_S, k^{(a)}_D\bigr) \sim \mathcal{N}\bigl(x^*,\ c\,H^{-1}\bigr),
\end{equation}
redrawing any sample with $k_S \le 0$. Two properties motivate this choice. First, patrons
vary widely along directions in which $J$ is flat and narrowly along directions in which
it is steep. Second, the spread has a direct meaning: if each patron's parameters were
adopted by a whole crowd,
\begin{equation}\label{eq:expected_loss}
  \mathbb{E}\bigl[J(x^{(a)}) - J^*\bigr]
  \approx \tfrac{1}{2}\operatorname{tr}\bigl(H \cdot c\,H^{-1}\bigr) = c,
\end{equation}
so $c$ is the average number of time steps lost per patron. Equivalently,
\cref{eq:population} is the quadratic approximation of a Boltzmann distribution
$p(x) \propto \exp\bigl(-(J(x) - J^*)/c\bigr)$, in which cheaper behaviors are more common
and $c$ acts as a temperature. We use $c = 1$. With the fitted
$x^* = (1.89, -3.02)$ and $H$, this gives % UPDATE
$\operatorname{sd}(k_S) \approx 0.35$, $\operatorname{sd}(k_D) \approx 2.45$ and
$\operatorname{corr}(k_S, k_D) \approx -0.76$: patrons agree closely on how strongly they
head for an exit but differ widely in how they respond to the crowd. % UPDATE

The fire weight must be positive, so we draw it from a lognormal distribution,
independently of $(k_S, k_D)$:
\begin{equation}\label{eq:kF}
  k^{(a)}_F = m\, e^{\sigma_F z_a}, \qquad z_a \sim \mathcal{N}(0, 1),
\end{equation}
with median $m$ and $\sigma_F = 0.5$, so about two thirds of patrons have $k_F$ within a
factor $e^{0.5} \approx 1.65$ of $m$. The median $m$ is learned in \cref{sec:learn_kF}.

\subsubsection{Learning the Fire Weight}\label{sec:learn_kF}

Consistent with \cref{as:consistent}, we hold the distribution of $(k_S, k_D)$ fixed and
learn only the median fire weight $m$. For one simulation, let $C$ be the casualty
fraction and $\bar t$ the mean exit time of survivors. We minimize
\begin{equation}\label{eq:fire_objective}
  \mathcal{J}(m) = \mathbb{E}\left[\, C + \lambda\, \frac{\bar t}{t_{\text{ref}}} \right],
  \qquad \lambda = 0.1,\ t_{\text{ref}} = 150 \text{ steps},
\end{equation}
where the expectation is over the ignition site, the population draw, the fire, and the
patrons' random choices; $t_{\text{ref}}$ is the optimal mean exit time without fire. With
these values one percent of casualties counts as much as $15$ steps of survivor exit time,
so casualties are the primary objective and exit time breaks near-ties.

We estimated $\mathcal{J}(m)$ for $m \in \{0, 0.5, 1, 2, 3, 5, 8, 12, 20\}$ with $R = 50$
replicates per value. Each replicate fixes an ignition site, a population draw (including
the standard normals $z_a$ in \cref{eq:kF}), a random stream for the patrons and a separate
random stream for the fire. These are shared across all values of $m$ (common random
numbers), and since the fire does not depend on the patrons, each replicate's fire evolves
identically for every $m$. Differences between values of $m$ are therefore paired
comparisons on identical scenarios.

\paragraph{Results.} \Cref{tab:kF} shows the results. % UPDATE all numbers below
Moderate fire avoidance roughly halves casualties: raising the median fire weight from
$m = 0$ to $m = 1$ lowers the casualty fraction from $2.3\%$ to $1.3\%$ (about $13$ of
$1010$ patrons). Strong avoidance is counterproductive: casualties rise to $3.1\%$ at
$m = 5$ and $3.8\%$ at $m = 20$, about three times the optimum, and the share of patrons
trapped with no route to an exit grows with $m$. Patrons who put distance from the fire
ahead of progress toward an exit are cut off or caught in worse positions. Fire avoidance
also delays survivors, whose mean exit time rises from $151$ steps at $m = 0$ to $180$ at
$m = 20$. $\mathcal{J}$ is minimized at $m = 1$, and no other value is within two paired
standard errors of it. Relative to $k_S^* \approx 1.9$, the best response weights a step
away from the fire about half as much as a step toward an exit.

\begin{table}[H]
  \centering
  \caption{Learning the median fire weight $m$ ($T_{\text{eng}} = \SI{6}{min}$, $50$
  replicates). Means over replicates; casualties are burned plus trapped, exit time is the
  survivors' mean in steps.}
  \label{tab:kF}
  % UPDATE: rerun learn_kf.py with the measured access points
  \begin{tabular}{rccccc}
    \toprule
    $m$ & $\mathcal{J}$ & casualties & burned & trapped & exit time \\
    \midrule
    0    & $0.124$ & $2.29\%$ & $1.93\%$ & $0.36\%$ & $151.1$ \\
    0.5  & $0.119$ & $1.66\%$ & $1.48\%$ & $0.18\%$ & $153.3$ \\
    1    & $0.116$ & $1.29\%$ & $1.08\%$ & $0.21\%$ & $155.0$ \\
    2    & $0.127$ & $1.87\%$ & $1.52\%$ & $0.35\%$ & $162.0$ \\
    3    & $0.135$ & $2.30\%$ & $1.88\%$ & $0.43\%$ & $168.0$ \\
    5    & $0.148$ & $3.09\%$ & $2.55\%$ & $0.54\%$ & $175.2$ \\
    8    & $0.153$ & $3.43\%$ & $2.84\%$ & $0.59\%$ & $177.7$ \\
    12   & $0.157$ & $3.79\%$ & $3.14\%$ & $0.65\%$ & $179.0$ \\
    20   & $0.159$ & $3.84\%$ & $3.16\%$ & $0.68\%$ & $180.1$ \\
    \bottomrule
  \end{tabular}
\end{table}

\paragraph{Sensitivity.} The optimum is robust to the fire model. With a faster fire
($T_{\text{eng}} = \SI{3}{min}$) and casualties occurring on first contact with the fire,
casualties are far higher ($13$--$16\%$), but $\mathcal{J}$ is again minimized at $m = 1$.
% UPDATE if rerun

\subsubsection{Constants}\label{sec:fire_constants}

\Cref{tab:fire_constants} lists every constant in the fire model and how it was determined.

\begin{table}[H]
  \centering
  \caption{Constants of the fire model.}
  \label{tab:fire_constants}
  \small
  \begin{tabular}{llll}
    \toprule
    Symbol & Meaning & Value & Source \\
    \midrule
    $\Delta t$ & time step & \SI{0.36}{s} & walking speed, cell size \\
    $\iota$ & ignite odds of fuel & $30$ & Minecraft (wool) \\
    $\beta$ & burn odds of fuel & $60$ & Minecraft (wool) \\
    $\kappa$ & difficulty level & $2$ & Minecraft (normal) \\
    $s$ & spread scale & $8$ & calibrated: $\approx 95\%$ of ignitions establish \\
    $T_{\text{eng}}$ & median engulfment time & \SI{6}{min} & assumed; \SI{3}{min} as sensitivity \\
    $q$ & fire update probability & $0.083$ & calibrated from $s$ and $T_{\text{eng}}$, \cref{eq:fire_q} \\
    $\ell$ & fire sensing range & $8$ cells (\SI{3.5}{m}) & assumed \\
    -- & exposure before casualty & \SI{5}{s} ($14$ steps) & assumed \\
    $T$ & time horizon & $2000$ steps & longer than any evacuation \\
    $x^*$ & center of $(k_S, k_D)$ & $(1.89, -3.02)$ & learned without fire % UPDATE
    \\
    $H$ & curvature of $J$ at $x^*$ & -- & fitted, \cref{eq:quadratic} \\
    $c$ & population spread & $1$ step & assumed \\
    $m$ & median fire weight & $1$ & learned, \cref{sec:learn_kF} % UPDATE
    \\
    $\sigma_F$ & spread of $k_F$ & $0.5$ & assumed \\
    $\lambda,\ t_{\text{ref}}$ & objective weights & $0.1,\ 150$ steps & chosen, \cref{eq:fire_objective} \\
    \bottomrule
  \end{tabular}
\end{table}

\subsubsection{Fire Assumptions}\label{sec:fire_assumptions}

\begin{assumption}
Every non-wall cell (seats, legroom, aisles, access points and exits) is equally
flammable, with the properties of wool in Minecraft; walls are stone and never burn.
\end{assumption}
\begin{assumption}
Fire spreads only between adjacent cells, and its spread does not depend on the patrons.
\end{assumption}
\begin{assumption}
The fire field represents untenable conditions (flames and smoke together), which spread
through the hall in a median of \SI{6}{min}.
\end{assumption}
\begin{assumption}
Patrons sense the fire within \SI{3.5}{m} and react to it in proportion to how many steps
away from it a move takes them.
\end{assumption}
\begin{assumption}
A patron who spends \SI{5}{s} in total in burning cells is incapacitated and becomes an
obstacle to others.
\end{assumption}
\begin{assumption}
Patrons' exit-seeking and crowd behavior, $(k_S, k_D)$, are the same as without fire
(\cref{as:consistent}); only their response to the fire, $k_F$, is new.
\end{assumption}

\subsubsection{Limitations}\label{sec:fire_limitations}

The fire model is two-dimensional and does not represent smoke, heat or the hot gas layer
under the ceiling, which cause most deaths in real assembly fires; the engulfment time is
a single stylized timescale for all of these. The Minecraft spread rules are not derived
from combustion physics, and their constants follow the game rather than measurements of
theater seating. Patrons respond only to fire within \SI{3.5}{m}; an alarm, the sight of
smoke across the hall or staff instructions would alert them sooner. Casualties are
binary, with no injuries or reduced walking speed. Finally, the optimal $(k_S, k_D)$ were
learned for a homogeneous crowd without fire, so the population in
\cref{eq:population} is near-optimal in that sense, not necessarily under fire.
 from main.tex. Needs: amsmath, siunitx, booktabs,
% cleveref, and the assumption environment defined in main.tex.
% Lines marked % UPDATE depend on the geometry and must be refreshed after
% rerunning learn.py, refine.py, calibrate_spread.py and learn_kf.py.
% ======================================================================

\subsection{Fire Scenario}\label{sec:fire}

We extend the evacuation model with a fire that starts at a random location in the
hall and spreads while patrons evacuate. The fire enters the model in three ways: it
makes cells impassable, it repels patrons through a fire proximity field, and it causes
casualties among patrons who remain in burning cells too long. Throughout, one time step
is $\Delta t = \SI{0.36}{s}$ and one cell is \SI{0.432}{m} square.

\subsubsection{Fire Spread}\label{sec:fire_spread}

We model fire with a stochastic cellular automaton adapted from the fire mechanics of
Minecraft (Java Edition) \cite{minecraftfire}, restricted to two dimensions. Each cell
$c$ carries an age $A_c \in \{-1, 0, 1, \dots, 15\}$, where $A_c = -1$ means $c$ is not
burning, and a flag $B_c \in \{0, 1\}$ recording whether its fuel has been used up. A cell
is \emph{flammable} if it is not a wall, not burning, and $B_c = 0$. Let $M(c)$ denote the
eight cells surrounding $c$.

At each time step the fire updates with probability $q$ (analogous to Minecraft's random
block ticks); when it updates, every cell changes simultaneously:
\begin{enumerate}
  \item \emph{Spread.} A burning cell $c'$ ignites each flammable neighbor independently
  with probability
  \begin{equation}\label{eq:fire_spread}
    p_{\text{ign}}(A_{c'}) = s \cdot \frac{\iota + 40 + 7\kappa}{100\,(A_{c'} + 30)},
  \end{equation}
  so a flammable cell $c$ catches fire with probability
  \begin{equation}\label{eq:fire_ignite}
    P_{\text{ign}}(c) = 1 - \prod_{\substack{c' \in M(c) \\ A_{c'} \ge 0}}
    \bigl(1 - p_{\text{ign}}(A_{c'})\bigr).
  \end{equation}
  \item \emph{Consumption.} A burning cell uses up its fuel with probability
  \begin{equation}\label{eq:fire_consume}
    p_{\text{cons}} = \frac{\beta}{300}.
  \end{equation}
  \item \emph{Aging.} The age of a burning cell increases by one with probability $1/3$,
  up to $15$.
  \item \emph{Burnout.} A burning cell goes out if $A_c > 3$ and no neighbor is flammable,
  or with probability $1/4$ once $A_c = 15$.
\end{enumerate}
A cell whose fire is consumed or burns out becomes ash: it is passable but no longer
flammable. Newly ignited cells start at age $0$. Each simulation begins with a single
ignition cell drawn uniformly at random from the non-wall, non-exit, unoccupied cells.

The functional forms of \cref{eq:fire_spread,eq:fire_consume} and the aging and burnout
rules follow Minecraft. Minecraft's fire is a block that sits next to separate fuel blocks
and can spread several blocks upward; in two dimensions we let the fire occupy the fuel
cell itself and spread within its eight-cell neighborhood. The constants are those of
wool in Minecraft, ignite odds $\iota = 30$ and burn odds $\beta = 60$, at the normal
difficulty level $\kappa = 2$ \cite{minecraftfire}, so a burning cell uses up its fuel
with probability $0.2$ per update. The spread scale $s$ and update probability $q$ are
calibrated in \cref{sec:fire_calibration}; Minecraft itself corresponds to $s = 1$.

\subsubsection{Calibrating the Fire}\label{sec:fire_calibration}

The spread scale $s$ and the update probability $q$ play different roles. The spread
scale controls the \emph{shape} of the fire: whether an ignition becomes an established
fire and how completely the fire front burns the hall. The update probability controls
only its \emph{speed}: lowering $q$ stretches the fire's timeline by a factor $1/q$
without changing how it spreads. We calibrate them separately.

\paragraph{Shape.} For each candidate $s$ we ran the fire alone from $40$ random ignition
sites, updating every step ($q = 1$). For each run we recorded the coverage $C(t)$, the
fraction of non-wall cells that have caught fire by update $t$, and its final value
$C_{\text{final}}$. We call a fire \emph{established} if $C_{\text{final}} \ge 0.5$, and
define its \emph{engulfment time} as the first $t$ with
$C(t) \ge 0.95\, C_{\text{final}}$. Because the front burns out behind itself and leaves
unburnt pockets, $C_{\text{final}} < 1$ even for fires that sweep the hall.
\Cref{tab:fire_calibration} shows the results. Below $s \approx 3$ almost no ignition
becomes established; the fire behaves like a branching process near its critical point,
so lowering $s$ makes fires fizzle rather than burn more slowly. We chose $s = 8$, for
which about $95\%$ of ignitions become established fires that burn essentially the whole
hall. This deliberately keeps a small share of fires that never take hold, as happens in
reality. Larger $s$ pushes $p_{\text{ign}}$ toward $1$, where the front becomes a
deterministic square determined by the grid rather than by the fire.

\begin{table}[H]
  \centering
  \caption{Fire calibration with the fire updating every step ($q = 1$), $40$ random
  ignitions per value of $s$. Coverage and engulfment updates are medians over
  established fires; $q$ is the update probability giving a median engulfment time of
  \SI{180}{s}.}
  \label{tab:fire_calibration}
  \begin{tabular}{rcccc}
    \toprule
    $s$ & established & coverage & engulfment (updates) & $q$ for \SI{180}{s} \\
    \midrule
    4  & $80\%$  & $0.90$ & $197$ & $0.394$ \\
    6  & $88\%$  & $0.99$ & $107$ & $0.214$ \\
    8  & $95\%$  & $1.00$ & $83$  & $0.166$ \\
    12 & $100\%$ & $1.00$ & $64$  & $0.128$ \\
    16 & $100\%$ & $1.00$ & $56$  & $0.113$ \\
    \bottomrule
  \end{tabular}
  % UPDATE: rerun calibrate_spread.py with the measured access points
\end{table}

\paragraph{Speed.} With $s = 8$ fixed, an established fire takes a median of $83$ updates
to engulf the hall, so an engulfment time $T_{\text{eng}}$ requires
\begin{equation}\label{eq:fire_q}
  q = \frac{83\,\Delta t}{T_{\text{eng}}} = 0.166 \cdot \frac{\SI{180}{s}}{T_{\text{eng}}}.
\end{equation}
We use $T_{\text{eng}} = \SI{6}{min}$, giving $q = 0.083$.

The engulfment time is the least certain input of the fire model, and we treat it as a
scenario parameter rather than a physical constant. For comparison, standard
fire-protection practice models fire growth as a $t^2$ design fire, with heat release
rate $Q(t) = \alpha t^2$; upholstered furniture is usually classed as a ``fast'' fire,
$\alpha = \SI{0.0469}{kW/s^2}$ \cite{sfpe}. If the burning area is $Q/q''$ for a heat
release rate per unit area $q''$, the flame front advances at
$v_f = \sqrt{\alpha/(\pi q'')}$, only a few millimeters per second for upholstered
seating. Flames that slow would not reach most patrons during an evacuation lasting about
a minute. In real assembly fires, smoke rather than flame spread is the main hazard, and
it fills a hall far faster than flames spread \cite{sfpe}. Our engulfment time should
therefore be read as the time for conditions throughout the hall to become untenable,
with the fire field standing in for flames and smoke together. We first used
$T_{\text{eng}} = \SI{3}{min}$, which produced $13$--$16\%$ casualties; we regard that as
implausibly severe and use \SI{6}{min}, reporting the \SI{3}{min} case as a sensitivity
check.

\subsubsection{Fire Proximity Field}\label{sec:fire_field}

Patrons should react to the fire before they reach it. Let $r_{ij}$ be the number of steps
from $(i, j)$ to the nearest burning cell through non-wall cells, using the eight-cell
neighborhood. The fire proximity field is
\begin{equation}\label{eq:fire_field}
  F_{ij} = \max\bigl(0,\ \ell - r_{ij}\bigr), \qquad \ell = 8 \text{ cells}
  \ (\approx \SI{3.5}{m}).
\end{equation}
Within $\ell$ cells of the fire, each step away from it lowers $F$ by exactly one, which is
the same unit in which each step toward an exit lowers $SD$. A fire weight $k_F$ is then
directly comparable to the exit weight $k_S$: $k_F / k_S$ is how much a patron values one
step away from the fire relative to one step toward an exit. An exponentially decaying
field would make $k_F$ depend on the distance at which it is evaluated and lose this
interpretation. Beyond $\ell$ the fire has no direct influence on movement.

\subsubsection{Movement and Casualties}\label{sec:fire_movement}

With fire present, the transition probability of patron $a$ at $(i, j)$ is
\begin{equation}\label{eq:trans_prob_fire}
  p^{(a)}_{i',j'} = N^{(a)}_{i,j}
  \exp\!\Bigl(k^{(a)}_D D_{i',j'} - k^{(a)}_S\, SD_{i',j'} - k^{(a)}_F F_{i',j'}\Bigr)
  \,\mathbf{1}\bigl[(i', j') \text{ admissible}\bigr],
\end{equation}
over the eight neighboring cells and the patron's own cell. A target cell is admissible if
it is walkable, unoccupied (by a patron or a casualty), and not burning. Because burning
cells and casualties block paths, $SD$ is recomputed whenever either changes. Cells from
which no exit can be reached are assigned one more than the largest finite distance, so
patrons cut off from every exit still move away from danger.

A patron standing in a burning cell accumulates one step of exposure per time step. When a
patron's cumulative exposure reaches \SI{5}{s} ($14$ steps), the patron becomes a
casualty and remains in the cell as an obstacle for the rest of the simulation; until
then, the patron can still walk out of the fire. Patrons still inside at the time horizon
$T = 2000$ steps (\SI{12}{min}) are counted as trapped. Both burned and trapped patrons
count as casualties.

\subsubsection{Heterogeneous Population}\label{sec:population}

Each patron $a$ receives individual weights $(k^{(a)}_S, k^{(a)}_D, k^{(a)}_F)$. By
\cref{as:optimal}, patrons behave approximately optimally under normal conditions, so we
draw $(k_S, k_D)$ from a neighborhood of the optimum $x^* = (k_S^*, k_D^*)$ for a
homogeneous crowd without fire. Near that optimum the mean exit time is well approximated
by a quadratic,
\begin{equation}\label{eq:quadratic}
  J(x) \approx J^* + \tfrac{1}{2}(x - x^*)^\top H\, (x - x^*),
\end{equation}
which we fit by least squares to the refined grid of simulations, obtaining $x^*$ and the
Hessian $H$. We draw
\begin{equation}\label{eq:population}
  \bigl(k^{(a)}_S, k^{(a)}_D\bigr) \sim \mathcal{N}\bigl(x^*,\ c\,H^{-1}\bigr),
\end{equation}
redrawing any sample with $k_S \le 0$. Two properties motivate this choice. First, patrons
vary widely along directions in which $J$ is flat and narrowly along directions in which
it is steep. Second, the spread has a direct meaning: if each patron's parameters were
adopted by a whole crowd,
\begin{equation}\label{eq:expected_loss}
  \mathbb{E}\bigl[J(x^{(a)}) - J^*\bigr]
  \approx \tfrac{1}{2}\operatorname{tr}\bigl(H \cdot c\,H^{-1}\bigr) = c,
\end{equation}
so $c$ is the average number of time steps lost per patron. Equivalently,
\cref{eq:population} is the quadratic approximation of a Boltzmann distribution
$p(x) \propto \exp\bigl(-(J(x) - J^*)/c\bigr)$, in which cheaper behaviors are more common
and $c$ acts as a temperature. We use $c = 1$. With the fitted
$x^* = (1.89, -3.02)$ and $H$, this gives % UPDATE
$\operatorname{sd}(k_S) \approx 0.35$, $\operatorname{sd}(k_D) \approx 2.45$ and
$\operatorname{corr}(k_S, k_D) \approx -0.76$: patrons agree closely on how strongly they
head for an exit but differ widely in how they respond to the crowd. % UPDATE

The fire weight must be positive, so we draw it from a lognormal distribution,
independently of $(k_S, k_D)$:
\begin{equation}\label{eq:kF}
  k^{(a)}_F = m\, e^{\sigma_F z_a}, \qquad z_a \sim \mathcal{N}(0, 1),
\end{equation}
with median $m$ and $\sigma_F = 0.5$, so about two thirds of patrons have $k_F$ within a
factor $e^{0.5} \approx 1.65$ of $m$. The median $m$ is learned in \cref{sec:learn_kF}.

\subsubsection{Learning the Fire Weight}\label{sec:learn_kF}

Consistent with \cref{as:consistent}, we hold the distribution of $(k_S, k_D)$ fixed and
learn only the median fire weight $m$. For one simulation, let $C$ be the casualty
fraction and $\bar t$ the mean exit time of survivors. We minimize
\begin{equation}\label{eq:fire_objective}
  \mathcal{J}(m) = \mathbb{E}\left[\, C + \lambda\, \frac{\bar t}{t_{\text{ref}}} \right],
  \qquad \lambda = 0.1,\ t_{\text{ref}} = 150 \text{ steps},
\end{equation}
where the expectation is over the ignition site, the population draw, the fire, and the
patrons' random choices; $t_{\text{ref}}$ is the optimal mean exit time without fire. With
these values one percent of casualties counts as much as $15$ steps of survivor exit time,
so casualties are the primary objective and exit time breaks near-ties.

We estimated $\mathcal{J}(m)$ for $m \in \{0, 0.5, 1, 2, 3, 5, 8, 12, 20\}$ with $R = 50$
replicates per value. Each replicate fixes an ignition site, a population draw (including
the standard normals $z_a$ in \cref{eq:kF}), a random stream for the patrons and a separate
random stream for the fire. These are shared across all values of $m$ (common random
numbers), and since the fire does not depend on the patrons, each replicate's fire evolves
identically for every $m$. Differences between values of $m$ are therefore paired
comparisons on identical scenarios.

\paragraph{Results.} \Cref{tab:kF} shows the results. % UPDATE all numbers below
Moderate fire avoidance roughly halves casualties: raising the median fire weight from
$m = 0$ to $m = 1$ lowers the casualty fraction from $2.3\%$ to $1.3\%$ (about $13$ of
$1010$ patrons). Strong avoidance is counterproductive: casualties rise to $3.1\%$ at
$m = 5$ and $3.8\%$ at $m = 20$, about three times the optimum, and the share of patrons
trapped with no route to an exit grows with $m$. Patrons who put distance from the fire
ahead of progress toward an exit are cut off or caught in worse positions. Fire avoidance
also delays survivors, whose mean exit time rises from $151$ steps at $m = 0$ to $180$ at
$m = 20$. $\mathcal{J}$ is minimized at $m = 1$, and no other value is within two paired
standard errors of it. Relative to $k_S^* \approx 1.9$, the best response weights a step
away from the fire about half as much as a step toward an exit.

\begin{table}[H]
  \centering
  \caption{Learning the median fire weight $m$ ($T_{\text{eng}} = \SI{6}{min}$, $50$
  replicates). Means over replicates; casualties are burned plus trapped, exit time is the
  survivors' mean in steps.}
  \label{tab:kF}
  % UPDATE: rerun learn_kf.py with the measured access points
  \begin{tabular}{rccccc}
    \toprule
    $m$ & $\mathcal{J}$ & casualties & burned & trapped & exit time \\
    \midrule
    0    & $0.124$ & $2.29\%$ & $1.93\%$ & $0.36\%$ & $151.1$ \\
    0.5  & $0.119$ & $1.66\%$ & $1.48\%$ & $0.18\%$ & $153.3$ \\
    1    & $0.116$ & $1.29\%$ & $1.08\%$ & $0.21\%$ & $155.0$ \\
    2    & $0.127$ & $1.87\%$ & $1.52\%$ & $0.35\%$ & $162.0$ \\
    3    & $0.135$ & $2.30\%$ & $1.88\%$ & $0.43\%$ & $168.0$ \\
    5    & $0.148$ & $3.09\%$ & $2.55\%$ & $0.54\%$ & $175.2$ \\
    8    & $0.153$ & $3.43\%$ & $2.84\%$ & $0.59\%$ & $177.7$ \\
    12   & $0.157$ & $3.79\%$ & $3.14\%$ & $0.65\%$ & $179.0$ \\
    20   & $0.159$ & $3.84\%$ & $3.16\%$ & $0.68\%$ & $180.1$ \\
    \bottomrule
  \end{tabular}
\end{table}

\paragraph{Sensitivity.} The optimum is robust to the fire model. With a faster fire
($T_{\text{eng}} = \SI{3}{min}$) and casualties occurring on first contact with the fire,
casualties are far higher ($13$--$16\%$), but $\mathcal{J}$ is again minimized at $m = 1$.
% UPDATE if rerun

\subsubsection{Constants}\label{sec:fire_constants}

\Cref{tab:fire_constants} lists every constant in the fire model and how it was determined.

\begin{table}[H]
  \centering
  \caption{Constants of the fire model.}
  \label{tab:fire_constants}
  \small
  \begin{tabular}{llll}
    \toprule
    Symbol & Meaning & Value & Source \\
    \midrule
    $\Delta t$ & time step & \SI{0.36}{s} & walking speed, cell size \\
    $\iota$ & ignite odds of fuel & $30$ & Minecraft (wool) \\
    $\beta$ & burn odds of fuel & $60$ & Minecraft (wool) \\
    $\kappa$ & difficulty level & $2$ & Minecraft (normal) \\
    $s$ & spread scale & $8$ & calibrated: $\approx 95\%$ of ignitions establish \\
    $T_{\text{eng}}$ & median engulfment time & \SI{6}{min} & assumed; \SI{3}{min} as sensitivity \\
    $q$ & fire update probability & $0.083$ & calibrated from $s$ and $T_{\text{eng}}$, \cref{eq:fire_q} \\
    $\ell$ & fire sensing range & $8$ cells (\SI{3.5}{m}) & assumed \\
    -- & exposure before casualty & \SI{5}{s} ($14$ steps) & assumed \\
    $T$ & time horizon & $2000$ steps & longer than any evacuation \\
    $x^*$ & center of $(k_S, k_D)$ & $(1.89, -3.02)$ & learned without fire % UPDATE
    \\
    $H$ & curvature of $J$ at $x^*$ & -- & fitted, \cref{eq:quadratic} \\
    $c$ & population spread & $1$ step & assumed \\
    $m$ & median fire weight & $1$ & learned, \cref{sec:learn_kF} % UPDATE
    \\
    $\sigma_F$ & spread of $k_F$ & $0.5$ & assumed \\
    $\lambda,\ t_{\text{ref}}$ & objective weights & $0.1,\ 150$ steps & chosen, \cref{eq:fire_objective} \\
    \bottomrule
  \end{tabular}
\end{table}

\subsubsection{Fire Assumptions}\label{sec:fire_assumptions}

\begin{assumption}
Every non-wall cell (seats, legroom, aisles, access points and exits) is equally
flammable, with the properties of wool in Minecraft; walls are stone and never burn.
\end{assumption}
\begin{assumption}
Fire spreads only between adjacent cells, and its spread does not depend on the patrons.
\end{assumption}
\begin{assumption}
The fire field represents untenable conditions (flames and smoke together), which spread
through the hall in a median of \SI{6}{min}.
\end{assumption}
\begin{assumption}
Patrons sense the fire within \SI{3.5}{m} and react to it in proportion to how many steps
away from it a move takes them.
\end{assumption}
\begin{assumption}
A patron who spends \SI{5}{s} in total in burning cells is incapacitated and becomes an
obstacle to others.
\end{assumption}
\begin{assumption}
Patrons' exit-seeking and crowd behavior, $(k_S, k_D)$, are the same as without fire
(\cref{as:consistent}); only their response to the fire, $k_F$, is new.
\end{assumption}

\subsubsection{Limitations}\label{sec:fire_limitations}

The fire model is two-dimensional and does not represent smoke, heat or the hot gas layer
under the ceiling, which cause most deaths in real assembly fires; the engulfment time is
a single stylized timescale for all of these. The Minecraft spread rules are not derived
from combustion physics, and their constants follow the game rather than measurements of
theater seating. Patrons respond only to fire within \SI{3.5}{m}; an alarm, the sight of
smoke across the hall or staff instructions would alert them sooner. Casualties are
binary, with no injuries or reduced walking speed. Finally, the optimal $(k_S, k_D)$ were
learned for a homogeneous crowd without fire, so the population in
\cref{eq:population} is near-optimal in that sense, not necessarily under fire.
 from main.tex. Needs: amsmath, siunitx, booktabs,
% cleveref, and the assumption environment defined in main.tex.
% Lines marked % UPDATE depend on the geometry and must be refreshed after
% rerunning learn.py, refine.py, calibrate_spread.py and learn_kf.py.
% ======================================================================

\subsection{Fire Scenario}\label{sec:fire}

We extend the evacuation model with a fire that starts at a random location in the
hall and spreads while patrons evacuate. The fire enters the model in three ways: it
makes cells impassable, it repels patrons through a fire proximity field, and it causes
casualties among patrons who remain in burning cells too long. Throughout, one time step
is $\Delta t = \SI{0.36}{s}$ and one cell is \SI{0.432}{m} square.

\subsubsection{Fire Spread}\label{sec:fire_spread}

We model fire with a stochastic cellular automaton adapted from the fire mechanics of
Minecraft (Java Edition) \cite{minecraftfire}, restricted to two dimensions. Each cell
$c$ carries an age $A_c \in \{-1, 0, 1, \dots, 15\}$, where $A_c = -1$ means $c$ is not
burning, and a flag $B_c \in \{0, 1\}$ recording whether its fuel has been used up. A cell
is \emph{flammable} if it is not a wall, not burning, and $B_c = 0$. Let $M(c)$ denote the
eight cells surrounding $c$.

At each time step the fire updates with probability $q$ (analogous to Minecraft's random
block ticks); when it updates, every cell changes simultaneously:
\begin{enumerate}
  \item \emph{Spread.} A burning cell $c'$ ignites each flammable neighbor independently
  with probability
  \begin{equation}\label{eq:fire_spread}
    p_{\text{ign}}(A_{c'}) = s \cdot \frac{\iota + 40 + 7\kappa}{100\,(A_{c'} + 30)},
  \end{equation}
  so a flammable cell $c$ catches fire with probability
  \begin{equation}\label{eq:fire_ignite}
    P_{\text{ign}}(c) = 1 - \prod_{\substack{c' \in M(c) \\ A_{c'} \ge 0}}
    \bigl(1 - p_{\text{ign}}(A_{c'})\bigr).
  \end{equation}
  \item \emph{Consumption.} A burning cell uses up its fuel with probability
  \begin{equation}\label{eq:fire_consume}
    p_{\text{cons}} = \frac{\beta}{300}.
  \end{equation}
  \item \emph{Aging.} The age of a burning cell increases by one with probability $1/3$,
  up to $15$.
  \item \emph{Burnout.} A burning cell goes out if $A_c > 3$ and no neighbor is flammable,
  or with probability $1/4$ once $A_c = 15$.
\end{enumerate}
A cell whose fire is consumed or burns out becomes ash: it is passable but no longer
flammable. Newly ignited cells start at age $0$. Each simulation begins with a single
ignition cell drawn uniformly at random from the non-wall, non-exit, unoccupied cells.

The functional forms of \cref{eq:fire_spread,eq:fire_consume} and the aging and burnout
rules follow Minecraft. Minecraft's fire is a block that sits next to separate fuel blocks
and can spread several blocks upward; in two dimensions we let the fire occupy the fuel
cell itself and spread within its eight-cell neighborhood. The constants are those of
wool in Minecraft, ignite odds $\iota = 30$ and burn odds $\beta = 60$, at the normal
difficulty level $\kappa = 2$ \cite{minecraftfire}, so a burning cell uses up its fuel
with probability $0.2$ per update. The spread scale $s$ and update probability $q$ are
calibrated in \cref{sec:fire_calibration}; Minecraft itself corresponds to $s = 1$.

\subsubsection{Calibrating the Fire}\label{sec:fire_calibration}

The spread scale $s$ and the update probability $q$ play different roles. The spread
scale controls the \emph{shape} of the fire: whether an ignition becomes an established
fire and how completely the fire front burns the hall. The update probability controls
only its \emph{speed}: lowering $q$ stretches the fire's timeline by a factor $1/q$
without changing how it spreads. We calibrate them separately.

\paragraph{Shape.} For each candidate $s$ we ran the fire alone from $40$ random ignition
sites, updating every step ($q = 1$). For each run we recorded the coverage $C(t)$, the
fraction of non-wall cells that have caught fire by update $t$, and its final value
$C_{\text{final}}$. We call a fire \emph{established} if $C_{\text{final}} \ge 0.5$, and
define its \emph{engulfment time} as the first $t$ with
$C(t) \ge 0.95\, C_{\text{final}}$. Because the front burns out behind itself and leaves
unburnt pockets, $C_{\text{final}} < 1$ even for fires that sweep the hall.
\Cref{tab:fire_calibration} shows the results. Below $s \approx 3$ almost no ignition
becomes established; the fire behaves like a branching process near its critical point,
so lowering $s$ makes fires fizzle rather than burn more slowly. We chose $s = 8$, for
which about $95\%$ of ignitions become established fires that burn essentially the whole
hall. This deliberately keeps a small share of fires that never take hold, as happens in
reality. Larger $s$ pushes $p_{\text{ign}}$ toward $1$, where the front becomes a
deterministic square determined by the grid rather than by the fire.

\begin{table}[H]
  \centering
  \caption{Fire calibration with the fire updating every step ($q = 1$), $40$ random
  ignitions per value of $s$. Coverage and engulfment updates are medians over
  established fires; $q$ is the update probability giving a median engulfment time of
  \SI{180}{s}.}
  \label{tab:fire_calibration}
  \begin{tabular}{rcccc}
    \toprule
    $s$ & established & coverage & engulfment (updates) & $q$ for \SI{180}{s} \\
    \midrule
    4  & $80\%$  & $0.90$ & $197$ & $0.394$ \\
    6  & $88\%$  & $0.99$ & $107$ & $0.214$ \\
    8  & $95\%$  & $1.00$ & $83$  & $0.166$ \\
    12 & $100\%$ & $1.00$ & $64$  & $0.128$ \\
    16 & $100\%$ & $1.00$ & $56$  & $0.113$ \\
    \bottomrule
  \end{tabular}
  % UPDATE: rerun calibrate_spread.py with the measured access points
\end{table}

\paragraph{Speed.} With $s = 8$ fixed, an established fire takes a median of $83$ updates
to engulf the hall, so an engulfment time $T_{\text{eng}}$ requires
\begin{equation}\label{eq:fire_q}
  q = \frac{83\,\Delta t}{T_{\text{eng}}} = 0.166 \cdot \frac{\SI{180}{s}}{T_{\text{eng}}}.
\end{equation}
We use $T_{\text{eng}} = \SI{6}{min}$, giving $q = 0.083$.

The engulfment time is the least certain input of the fire model, and we treat it as a
scenario parameter rather than a physical constant. For comparison, standard
fire-protection practice models fire growth as a $t^2$ design fire, with heat release
rate $Q(t) = \alpha t^2$; upholstered furniture is usually classed as a ``fast'' fire,
$\alpha = \SI{0.0469}{kW/s^2}$ \cite{sfpe}. If the burning area is $Q/q''$ for a heat
release rate per unit area $q''$, the flame front advances at
$v_f = \sqrt{\alpha/(\pi q'')}$, only a few millimeters per second for upholstered
seating. Flames that slow would not reach most patrons during an evacuation lasting about
a minute. In real assembly fires, smoke rather than flame spread is the main hazard, and
it fills a hall far faster than flames spread \cite{sfpe}. Our engulfment time should
therefore be read as the time for conditions throughout the hall to become untenable,
with the fire field standing in for flames and smoke together. We first used
$T_{\text{eng}} = \SI{3}{min}$, which produced $13$--$16\%$ casualties; we regard that as
implausibly severe and use \SI{6}{min}, reporting the \SI{3}{min} case as a sensitivity
check.

\subsubsection{Fire Proximity Field}\label{sec:fire_field}

Patrons should react to the fire before they reach it. Let $r_{ij}$ be the number of steps
from $(i, j)$ to the nearest burning cell through non-wall cells, using the eight-cell
neighborhood. The fire proximity field is
\begin{equation}\label{eq:fire_field}
  F_{ij} = \max\bigl(0,\ \ell - r_{ij}\bigr), \qquad \ell = 8 \text{ cells}
  \ (\approx \SI{3.5}{m}).
\end{equation}
Within $\ell$ cells of the fire, each step away from it lowers $F$ by exactly one, which is
the same unit in which each step toward an exit lowers $SD$. A fire weight $k_F$ is then
directly comparable to the exit weight $k_S$: $k_F / k_S$ is how much a patron values one
step away from the fire relative to one step toward an exit. An exponentially decaying
field would make $k_F$ depend on the distance at which it is evaluated and lose this
interpretation. Beyond $\ell$ the fire has no direct influence on movement.

\subsubsection{Movement and Casualties}\label{sec:fire_movement}

With fire present, the transition probability of patron $a$ at $(i, j)$ is
\begin{equation}\label{eq:trans_prob_fire}
  p^{(a)}_{i',j'} = N^{(a)}_{i,j}
  \exp\!\Bigl(k^{(a)}_D D_{i',j'} - k^{(a)}_S\, SD_{i',j'} - k^{(a)}_F F_{i',j'}\Bigr)
  \,\mathbf{1}\bigl[(i', j') \text{ admissible}\bigr],
\end{equation}
over the eight neighboring cells and the patron's own cell. A target cell is admissible if
it is walkable, unoccupied (by a patron or a casualty), and not burning. Because burning
cells and casualties block paths, $SD$ is recomputed whenever either changes. Cells from
which no exit can be reached are assigned one more than the largest finite distance, so
patrons cut off from every exit still move away from danger.

A patron standing in a burning cell accumulates one step of exposure per time step. When a
patron's cumulative exposure reaches \SI{5}{s} ($14$ steps), the patron becomes a
casualty and remains in the cell as an obstacle for the rest of the simulation; until
then, the patron can still walk out of the fire. Patrons still inside at the time horizon
$T = 2000$ steps (\SI{12}{min}) are counted as trapped. Both burned and trapped patrons
count as casualties.

\subsubsection{Heterogeneous Population}\label{sec:population}

Each patron $a$ receives individual weights $(k^{(a)}_S, k^{(a)}_D, k^{(a)}_F)$. By
\cref{as:optimal}, patrons behave approximately optimally under normal conditions, so we
draw $(k_S, k_D)$ from a neighborhood of the optimum $x^* = (k_S^*, k_D^*)$ for a
homogeneous crowd without fire. Near that optimum the mean exit time is well approximated
by a quadratic,
\begin{equation}\label{eq:quadratic}
  J(x) \approx J^* + \tfrac{1}{2}(x - x^*)^\top H\, (x - x^*),
\end{equation}
which we fit by least squares to the refined grid of simulations, obtaining $x^*$ and the
Hessian $H$. We draw
\begin{equation}\label{eq:population}
  \bigl(k^{(a)}_S, k^{(a)}_D\bigr) \sim \mathcal{N}\bigl(x^*,\ c\,H^{-1}\bigr),
\end{equation}
redrawing any sample with $k_S \le 0$. Two properties motivate this choice. First, patrons
vary widely along directions in which $J$ is flat and narrowly along directions in which
it is steep. Second, the spread has a direct meaning: if each patron's parameters were
adopted by a whole crowd,
\begin{equation}\label{eq:expected_loss}
  \mathbb{E}\bigl[J(x^{(a)}) - J^*\bigr]
  \approx \tfrac{1}{2}\operatorname{tr}\bigl(H \cdot c\,H^{-1}\bigr) = c,
\end{equation}
so $c$ is the average number of time steps lost per patron. Equivalently,
\cref{eq:population} is the quadratic approximation of a Boltzmann distribution
$p(x) \propto \exp\bigl(-(J(x) - J^*)/c\bigr)$, in which cheaper behaviors are more common
and $c$ acts as a temperature. We use $c = 1$. With the fitted
$x^* = (1.89, -3.02)$ and $H$, this gives % UPDATE
$\operatorname{sd}(k_S) \approx 0.35$, $\operatorname{sd}(k_D) \approx 2.45$ and
$\operatorname{corr}(k_S, k_D) \approx -0.76$: patrons agree closely on how strongly they
head for an exit but differ widely in how they respond to the crowd. % UPDATE

The fire weight must be positive, so we draw it from a lognormal distribution,
independently of $(k_S, k_D)$:
\begin{equation}\label{eq:kF}
  k^{(a)}_F = m\, e^{\sigma_F z_a}, \qquad z_a \sim \mathcal{N}(0, 1),
\end{equation}
with median $m$ and $\sigma_F = 0.5$, so about two thirds of patrons have $k_F$ within a
factor $e^{0.5} \approx 1.65$ of $m$. The median $m$ is learned in \cref{sec:learn_kF}.

\subsubsection{Learning the Fire Weight}\label{sec:learn_kF}

Consistent with \cref{as:consistent}, we hold the distribution of $(k_S, k_D)$ fixed and
learn only the median fire weight $m$. For one simulation, let $C$ be the casualty
fraction and $\bar t$ the mean exit time of survivors. We minimize
\begin{equation}\label{eq:fire_objective}
  \mathcal{J}(m) = \mathbb{E}\left[\, C + \lambda\, \frac{\bar t}{t_{\text{ref}}} \right],
  \qquad \lambda = 0.1,\ t_{\text{ref}} = 150 \text{ steps},
\end{equation}
where the expectation is over the ignition site, the population draw, the fire, and the
patrons' random choices; $t_{\text{ref}}$ is the optimal mean exit time without fire. With
these values one percent of casualties counts as much as $15$ steps of survivor exit time,
so casualties are the primary objective and exit time breaks near-ties.

We estimated $\mathcal{J}(m)$ for $m \in \{0, 0.5, 1, 2, 3, 5, 8, 12, 20\}$ with $R = 50$
replicates per value. Each replicate fixes an ignition site, a population draw (including
the standard normals $z_a$ in \cref{eq:kF}), a random stream for the patrons and a separate
random stream for the fire. These are shared across all values of $m$ (common random
numbers), and since the fire does not depend on the patrons, each replicate's fire evolves
identically for every $m$. Differences between values of $m$ are therefore paired
comparisons on identical scenarios.

\paragraph{Results.} \Cref{tab:kF} shows the results. % UPDATE all numbers below
Moderate fire avoidance roughly halves casualties: raising the median fire weight from
$m = 0$ to $m = 1$ lowers the casualty fraction from $2.3\%$ to $1.3\%$ (about $13$ of
$1010$ patrons). Strong avoidance is counterproductive: casualties rise to $3.1\%$ at
$m = 5$ and $3.8\%$ at $m = 20$, about three times the optimum, and the share of patrons
trapped with no route to an exit grows with $m$. Patrons who put distance from the fire
ahead of progress toward an exit are cut off or caught in worse positions. Fire avoidance
also delays survivors, whose mean exit time rises from $151$ steps at $m = 0$ to $180$ at
$m = 20$. $\mathcal{J}$ is minimized at $m = 1$, and no other value is within two paired
standard errors of it. Relative to $k_S^* \approx 1.9$, the best response weights a step
away from the fire about half as much as a step toward an exit.

\begin{table}[H]
  \centering
  \caption{Learning the median fire weight $m$ ($T_{\text{eng}} = \SI{6}{min}$, $50$
  replicates). Means over replicates; casualties are burned plus trapped, exit time is the
  survivors' mean in steps.}
  \label{tab:kF}
  % UPDATE: rerun learn_kf.py with the measured access points
  \begin{tabular}{rccccc}
    \toprule
    $m$ & $\mathcal{J}$ & casualties & burned & trapped & exit time \\
    \midrule
    0    & $0.124$ & $2.29\%$ & $1.93\%$ & $0.36\%$ & $151.1$ \\
    0.5  & $0.119$ & $1.66\%$ & $1.48\%$ & $0.18\%$ & $153.3$ \\
    1    & $0.116$ & $1.29\%$ & $1.08\%$ & $0.21\%$ & $155.0$ \\
    2    & $0.127$ & $1.87\%$ & $1.52\%$ & $0.35\%$ & $162.0$ \\
    3    & $0.135$ & $2.30\%$ & $1.88\%$ & $0.43\%$ & $168.0$ \\
    5    & $0.148$ & $3.09\%$ & $2.55\%$ & $0.54\%$ & $175.2$ \\
    8    & $0.153$ & $3.43\%$ & $2.84\%$ & $0.59\%$ & $177.7$ \\
    12   & $0.157$ & $3.79\%$ & $3.14\%$ & $0.65\%$ & $179.0$ \\
    20   & $0.159$ & $3.84\%$ & $3.16\%$ & $0.68\%$ & $180.1$ \\
    \bottomrule
  \end{tabular}
\end{table}

\paragraph{Sensitivity.} The optimum is robust to the fire model. With a faster fire
($T_{\text{eng}} = \SI{3}{min}$) and casualties occurring on first contact with the fire,
casualties are far higher ($13$--$16\%$), but $\mathcal{J}$ is again minimized at $m = 1$.
% UPDATE if rerun

\subsubsection{Constants}\label{sec:fire_constants}

\Cref{tab:fire_constants} lists every constant in the fire model and how it was determined.

\begin{table}[H]
  \centering
  \caption{Constants of the fire model.}
  \label{tab:fire_constants}
  \small
  \begin{tabular}{llll}
    \toprule
    Symbol & Meaning & Value & Source \\
    \midrule
    $\Delta t$ & time step & \SI{0.36}{s} & walking speed, cell size \\
    $\iota$ & ignite odds of fuel & $30$ & Minecraft (wool) \\
    $\beta$ & burn odds of fuel & $60$ & Minecraft (wool) \\
    $\kappa$ & difficulty level & $2$ & Minecraft (normal) \\
    $s$ & spread scale & $8$ & calibrated: $\approx 95\%$ of ignitions establish \\
    $T_{\text{eng}}$ & median engulfment time & \SI{6}{min} & assumed; \SI{3}{min} as sensitivity \\
    $q$ & fire update probability & $0.083$ & calibrated from $s$ and $T_{\text{eng}}$, \cref{eq:fire_q} \\
    $\ell$ & fire sensing range & $8$ cells (\SI{3.5}{m}) & assumed \\
    -- & exposure before casualty & \SI{5}{s} ($14$ steps) & assumed \\
    $T$ & time horizon & $2000$ steps & longer than any evacuation \\
    $x^*$ & center of $(k_S, k_D)$ & $(1.89, -3.02)$ & learned without fire % UPDATE
    \\
    $H$ & curvature of $J$ at $x^*$ & -- & fitted, \cref{eq:quadratic} \\
    $c$ & population spread & $1$ step & assumed \\
    $m$ & median fire weight & $1$ & learned, \cref{sec:learn_kF} % UPDATE
    \\
    $\sigma_F$ & spread of $k_F$ & $0.5$ & assumed \\
    $\lambda,\ t_{\text{ref}}$ & objective weights & $0.1,\ 150$ steps & chosen, \cref{eq:fire_objective} \\
    \bottomrule
  \end{tabular}
\end{table}

\subsubsection{Fire Assumptions}\label{sec:fire_assumptions}

\begin{assumption}
Every non-wall cell (seats, legroom, aisles, access points and exits) is equally
flammable, with the properties of wool in Minecraft; walls are stone and never burn.
\end{assumption}
\begin{assumption}
Fire spreads only between adjacent cells, and its spread does not depend on the patrons.
\end{assumption}
\begin{assumption}
The fire field represents untenable conditions (flames and smoke together), which spread
through the hall in a median of \SI{6}{min}.
\end{assumption}
\begin{assumption}
Patrons sense the fire within \SI{3.5}{m} and react to it in proportion to how many steps
away from it a move takes them.
\end{assumption}
\begin{assumption}
A patron who spends \SI{5}{s} in total in burning cells is incapacitated and becomes an
obstacle to others.
\end{assumption}
\begin{assumption}
Patrons' exit-seeking and crowd behavior, $(k_S, k_D)$, are the same as without fire
(\cref{as:consistent}); only their response to the fire, $k_F$, is new.
\end{assumption}

\subsubsection{Limitations}\label{sec:fire_limitations}

The fire model is two-dimensional and does not represent smoke, heat or the hot gas layer
under the ceiling, which cause most deaths in real assembly fires; the engulfment time is
a single stylized timescale for all of these. The Minecraft spread rules are not derived
from combustion physics, and their constants follow the game rather than measurements of
theater seating. Patrons respond only to fire within \SI{3.5}{m}; an alarm, the sight of
smoke across the hall or staff instructions would alert them sooner. Casualties are
binary, with no injuries or reduced walking speed. Finally, the optimal $(k_S, k_D)$ were
learned for a homogeneous crowd without fire, so the population in
\cref{eq:population} is near-optimal in that sense, not necessarily under fire.
